\section{What the Gates Caught: A Failure Catalogue}\label{sec:failures}

Every controlled result in this paper passed through the same mechanical gates: implausibility checks on scored verdicts, a fixed held-out evaluation set required before any win call, iso-FLOP confound checks, and a claim-to-file audit that opens every cited artifact. This section catalogues what those gates caught in the project's own work. Each entry gives the symptom, the detection mechanism, the root cause, the fix, and the guard that remains.

\paragraph{1. Token-level shuffle in the first data-mix run.}
Symptom: the 50/50 dclm-edu/FineWeb-Edu mix arm scored roughly $2\times$ \emph{worse} than both pure parents (wikitext-2 BPB 2.77 against $\approx$1.52, as recorded in the fix commit) -- physically implausible for a blend, whose loss should sit near the convex hull of its parents. Detection: the implausibility check on the verdict flagged it before any conclusion was recorded. Root cause: a \texttt{numpy} shuffle applied at the \emph{token} level interleaved the two sources token by token, destroying all sequence structure. Fix: commit \texttt{a576baa} concatenates the two coherent halves and mixes at the packed-sequence level, with the loader shuffling whole 4096-token windows. A 30-step probe after the fix shows the training cross-entropy descending from 11.65 to 8.04, where the token-soup variant had stayed flat near 11. The broken cohort was discarded and rerun; the corrected mix result appears in Table~\ref{tab:data}. Guard: the loader now carries an in-code warning, and the implausibility check stays in the scoring path.

\paragraph{2. The SFT eval-token confound.}
Symptom: in-loop advisory scoring showed response-masked SFT beating its iso-FLOP control by 0.68 perplexity, nominally significant (Figure~\ref{fig:sft-confound}; in-loop values, not suite-stamped). Detection: the protocol forbids a win call from in-loop numbers; a fixed held-out set must be scored first. Root cause: the SFT arm was scored on response tokens only while the control was scored on all tokens -- two incomparable token populations. On one fixed held-out response-masked set the gap collapses to $+0.009$ perplexity (95\% CI $[+0.0045, +0.0139]$, 3 seeds, paired), and a full-sequence cross-check gives $-0.006$ in the control's favor, not significant; the two disagree, so the verdict of record is directional (Section~\ref{sec:sft}, Table~\ref{tab:sft}). Guard: in-loop perplexities are permanently labeled advisory.

\paragraph{3. The RFT smoke-log misread.}
Symptom: an earlier draft of the RLVR verdict stated that the rejection-fine-tuning control collected zero verifier-correct completions. Detection: the claim-to-file audit for this paper re-read the full training log rather than its opening lines. Root cause: the ``0 collected'' line belongs to the one-step smoke test of 2026-07-02 at the top of the same log; the real 300-step run collected 352 of 38{,}400 rollouts ($\approx$0.9\%, matching the flat $\approx$0.9\% GRPO training reward; Table~\ref{tab:rlvr}). Fix: the verdict file and the repository README were corrected on 2026-07-06; the correction note is embedded in the regenerated artifact. The null verdict is unchanged -- 352 completions were too few to separate RFT from the floor -- but ``zero'' was wrong. Guard: smoke-test output is now read as a distinct log segment.

\paragraph{4. The README data-ratio pairing error.}
Symptom: repository documentation paired the $2.14\times$ gap-to-original with ``$\approx$275{,}000$\times$ less data.'' Detection: the claim-to-file audit; the source file pairs $275{,}000\times$ with the 131M-token Phase-A run, whose gap is $3.5\times$. Root cause: numbers from two different runs merged into one sentence. The correct pairing is $2.14\times$ at $\approx$30{,}000$\times$ less data (1.19B tokens) and $3.5\times$ at $\approx$275{,}000$\times$ less data (131M tokens), as in Figure~\ref{fig:repro-gap}; both gaps are ratios of in-loop validation perplexity, single run each, and are descriptive only. Corrected 2026-07-06.

\paragraph{5. The incomplete partial-RoPE-0.10 arm.}
The exploratory partial-RoPE-0.10 build died at step 5{,}450 of 18{,}150 ($\approx$30\% of budget), with a last-eval in-loop validation perplexity of 50.71 (single run, not suite-stamped). It is reported as incomplete in Section~\ref{sec:pretrain} rather than silently dropped -- an honesty entry rather than a gate catch.

\paragraph{6. A stale derived field in the ledger.}
Symptom: the ledger entry for the mid-training anneal carries a sub-field \texttt{c25\_cap} reading ``directional'' beside a verdict field reading ``win.'' Detection: the audit cross-checked every ledger field against the run's verdict artifact. Root cause: the cap predates the same-night upgrade that added the effective-context-length ladder and made the stage's required evaluation battery complete; the verdict fields were updated, the derived sub-field was not. The authoritative verdict file (no missing hard items) supports the win reported in Table~\ref{tab:midtrain}. Guard: noted as a standing hazard -- derived fields duplicated across artifacts can go stale independently of the record they summarize.

\paragraph{7. A parameter-count citation pointing at the wrong file.}
Symptom: the model README cited \texttt{verify.json} as evidence for the 596{,}049{,}920-parameter count (Table~\ref{tab:repro}), but that file contains no parameter field. The number itself is correct -- it is recovered exactly by recomputation from the architecture configuration and asserted in the build's test suite -- only the citation was off. Detection: the claim-to-file audit. Guard: every number in this paper traces to the file that actually contains it.

None of these seven was caught by re-reading prose. Each was caught by a mechanical gate: an implausibility check on a scored verdict, a fixed held-out set required before any win call, or a claim-to-file audit that opens every cited artifact. These are the same gates that validate the positive results of Sections \ref{sec:pretrain}--\ref{sec:rlvr}; a protocol that demonstrably catches its own errors is the strongest warrant this study can offer for the claims that survived it.
